% quiz_answers.tex -- answers and worked solutions for quiz.tex, same layout (quiz.sty).
% The worked answers are the same text as src/m5t3_booklet/answers.tex.
% Every number: checks/m5t3-booklet/numbers.py
% Build: TEXINPUTS=.: xelatex quiz_answers.tex (twice)
\documentclass[12pt]{article}
\usepackage{quiz}
\hypersetup{pdftitle={Multiple Choice Quiz - Module 5 Topic 3 - Gravity - Answers}}
% links between the two files are not possible, so question headings are plain text here
\renewcommand\hyperlink[2]{#2}

\begin{document}
\quiztitle{Year 12 Physics}{Multiple Choice Quiz - Module 5 Topic 3 - Gravity: Answers}

\begin{center}
  \renewcommand\arraystretch{1.5}
  \begin{tabular}{|l|*{10}{>{\centering\arraybackslash}p{7.5mm}|}}
    \hline
    \textbf{Question} & 1 & 2 & 3 & 4 & 5 & 6 & 7 & 8 & 9 & 10\\ \hline
    \textbf{Answer} & C & D & A & B & D & C & A & D & D & B\\ \hline\hline
    \textbf{Question} & 11 & 12 & 13 & 14 & 15 & 16 & 17 & 18 & 19 & 20\\ \hline
    \textbf{Answer} & C & D & A & C & A & B & C & D & A & B\\ \hline\hline
    \textbf{Question} & 21 & 22 & 23 & 24 & 25 & 26 & 27 & 28 & 29 & 30\\ \hline
    \textbf{Answer} & C & D & B & C & A & C & A & B & D & B\\ \hline
  \end{tabular}
\end{center}
\vspace{4mm}

\begin{practiceanswers}[Worked answers]

% ===================================================================== Section A
% ---- Q1
\Needspace{9\baselineskip}
\begin{practiceans}
  \correct{C}
  Newton's third law: Earth pulls on the satellite, and the satellite pulls on Earth. These two
  forces are equal in size and opposite in direction. But $a = \dfrac{F}{m}$. Earth's mass
  (\qty{6.0e24}{kg}) is huge, so the same force gives Earth a tiny acceleration.
  \begin{itemize}
    \item \xmark\ A: the two forces of a third-law pair are \emph{always} equal in size.
    \item \xmark\ B: every mass pulls on every other mass. $F = \dfrac{GMm}{r^2}$ is not zero.
    \item \xmark\ D: equal forces do not give equal accelerations, because the masses are very
          different.
  \end{itemize}
\end{practiceans}

% ---- Q2
\Needspace{9\baselineskip}
\begin{practiceans}
  \correct{D}
  $g = \dfrac{GM}{r^2}$. Mass $\times 2$ makes $g$ 2 times bigger. Radius $\times \frac{1}{2}$
  makes $r^2$ 4 times smaller, so $g$ is 4 times bigger. Together: $2 \times 4 = 8$, so the field
  strength is $8g$.
  \begin{itemize}
    \item \xmark\ A ($\frac{g}{2}$): multiplies by $\left(\frac{1}{2}\right)^2$ instead of dividing by it.
    \item \xmark\ B ($2g$): ignores the change in radius.
    \item \xmark\ C ($4g$): divides by $r$, not by $r^2$: $2 \div \frac{1}{2} = 4$.
  \end{itemize}
\end{practiceans}

% ---- Q3
\Needspace{9\baselineskip}
\begin{practiceans}
  \correct{A}
  $g = \dfrac{GM}{r^2}$, so $g \propto \dfrac{1}{r^2}$. For $\dfrac{g}{9}$, $r^2$ must be 9 times
  bigger, so $r$ is 3 times bigger: $r = 3r_\mathrm{E}$ from the \emph{centre} of Earth. The
  altitude is measured from the surface: $3r_\mathrm{E} - r_\mathrm{E} = 2r_\mathrm{E}$.
  \begin{itemize}
    \item \xmark\ B ($3r_\mathrm{E}$): this is the distance from the centre, not the altitude.
    \item \xmark\ C ($8r_\mathrm{E}$): uses $g \propto \frac{1}{r}$ (forgets the square), then subtracts
          $r_\mathrm{E}$.
    \item \xmark\ D ($9r_\mathrm{E}$): uses $g \propto \frac{1}{r}$ and does not subtract $r_\mathrm{E}$.
  \end{itemize}
\end{practiceans}

% ---- Q4
\Needspace{9\baselineskip}
\begin{practiceans}
  \correct{B}
  \[ F = \frac{GMm}{r^2}
       = \frac{(\num{6.67e-11})(\num{6.42e23})(\num{1.07e16})}{(\num{9.38e6})^2}
       = \qty{5.2e15}{N} \]
  The period is extra data. You can use it as a check: gravity gives the centripetal force, and
  $\dfrac{4\pi^2 m r}{T^2}$ also gives about \qty{5.2e15}{N}.
  \begin{itemize}
    \item \xmark\ A (\qty{0.49}{N}): this is $\dfrac{GM}{r^2}$, the field strength in
          \si{N.kg^{-1}}. It forgets to multiply by the mass of Phobos.
    \item \xmark\ C: uses $r$ in kilometres instead of metres.
    \item \xmark\ D: forgets to square $r$.
  \end{itemize}
\end{practiceans}

% ---- Q5
\Needspace{9\baselineskip}
\begin{practiceans}
  \correct{D}
  Square both sides of the formula:
  \[ T^2 = \frac{4\pi^2}{g}\, l \]
  This has the form $y = kx$. So $T^2$ against $l$ is a straight line through the origin with
  gradient $\dfrac{4\pi^2}{g}$. On the Moon $g$ is smaller, so the gradient is bigger (about 6
  times). The Moon line (dashed) is the steeper one.
  \begin{itemize}
    \item \xmark\ A: straight lines, but the Earth line is steeper. That would need a smaller $g$ on
          Earth.
    \item \xmark\ B and C: these curves have the shape of $T$ against $l$ ($T \propto \sqrt{l}$),
          not $T^2$ against $l$.
  \end{itemize}
\end{practiceans}

% ---- Q6
\Needspace{9\baselineskip}
\begin{practiceans}
  \correct{C}
  At this point, the pull of Earth equals the pull of the Moon:
  \[ \frac{GM_\mathrm{E}m}{d_\mathrm{E}^2} = \frac{GM_\mathrm{M}m}{d_\mathrm{M}^2}
     \quad\Rightarrow\quad \frac{d_\mathrm{E}^2}{d_\mathrm{M}^2} = \frac{M_\mathrm{E}}{M_\mathrm{M}} = 81
     \quad\Rightarrow\quad d_\mathrm{E} = 9\,d_\mathrm{M} \]
  The whole distance is $d_\mathrm{E} + d_\mathrm{M} = 10\,d_\mathrm{M}$. So the point is
  $\frac{9}{10}$ of the way from Earth. It is close to the Moon, because the Moon's pull is weaker.
  \begin{itemize}
    \item \xmark\ A: halfway is correct only if the two masses are equal.
    \item \xmark\ B: close to Earth, Earth's pull is much stronger. $\frac{1}{10}$ is the distance
          from the \emph{Moon}.
    \item \xmark\ D: uses $F \propto \frac{1}{r}$, so it forgets the square root.
  \end{itemize}
\end{practiceans}

% ===================================================================== Section B
% ---- Q7
\Needspace{9\baselineskip}
\begin{practiceans}
  \correct{A}
  Only one real force acts on the satellite: gravity, towards the centre of Earth. Gravity
  \emph{is} the centripetal force. “Centripetal force” is just the name for the net force
  towards the centre. It is not an extra force. No force is needed along the direction of motion:
  the satellite keeps moving because of its inertia.
  \begin{itemize}
    \item \xmark\ B: there is no outward force. If the forces balanced, the satellite would move in a
          straight line.
    \item \xmark\ C: in a stable orbit there is no air resistance, so no thrust is needed.
    \item \xmark\ D: $F_\mathrm{c}$ and $F_\mathrm{g}$ are the same force. Drawing both counts gravity
          twice.
  \end{itemize}
\end{practiceans}

% ---- Q8
\Needspace{9\baselineskip}
\begin{practiceans}
  \correct{D}
  A geostationary satellite stays above the same point on the equator. To do this, it must turn
  with Earth: same direction, same period (about 24 hours), in an orbit above the equator. Its
  altitude is about \qty{36000}{km}.
  \begin{itemize}
    \item \xmark\ A: the plane of every orbit passes through the centre of Earth. Only an orbit
          above the equator can keep a satellite above the same point. Sydney is not on the
          equator.
    \item \xmark\ B: low Earth orbit is about 200 to \qty{2000}{km} high. Geostationary orbit is
          much higher.
    \item \xmark\ C: $v = \sqrt{\dfrac{GM}{r}}$, so a higher orbit is slower. GEO is about
          \qty{3.1}{km.s^{-1}}, LEO about \qty{7.7}{km.s^{-1}}.
  \end{itemize}
\end{practiceans}

% ---- Q9
\Needspace{9\baselineskip}
\begin{practiceans}
  \correct{D}
  $v = \sqrt{\dfrac{GM}{r}}$. The speed depends only on the planet's mass $M$ and the radius $r$,
  not on the satellite's mass. Same orbit, so same speed (and same period).
  \begin{itemize}
    \item \xmark\ A: $K = \frac{1}{2}mv^2$, so the heavier satellite has 3 times the kinetic energy.
    \item \xmark\ B: $F = \dfrac{GMm}{r^2}$, so the force on the heavier satellite is 3 times bigger.
    \item \xmark\ C: $U = -\dfrac{GMm}{r}$, so the heavier satellite's potential energy is 3 times bigger (more negative).
  \end{itemize}
\end{practiceans}

% ---- Q10
\Needspace{9\baselineskip}
\begin{practiceans}
  \correct{B}
  $v = \sqrt{\dfrac{GM}{r}}$, so $v \propto \dfrac{1}{\sqrt{r}}$. The radius becomes 9 times
  bigger, so the speed is divided by $\sqrt{9} = 3$. The new speed is $\dfrac{v}{3}$.
  \begin{itemize}
    \item \xmark\ A: forgets the square root.
    \item \xmark\ C and D: a higher orbit is \emph{slower}, not faster.
  \end{itemize}
\end{practiceans}

% ---- Q11
\Needspace{9\baselineskip}
\begin{practiceans}
  \correct{C}
  $\dfrac{r^3}{T^2}$ is the same for all moons of one planet, so $T^2 \propto r^3$. Radius
  $\times 4$ gives $T^2 \times 4^3 = 64$, so $T \times \sqrt{64} = 8$. The period is $8T$.
  \begin{itemize}
    \item \xmark\ A ($2T$): uses $T \propto \sqrt{r}$.
    \item \xmark\ B ($4T$): uses $T \propto r$.
    \item \xmark\ D ($16T$): uses $T \propto r^2$.
  \end{itemize}
\end{practiceans}

% ---- Q12
\Needspace{9\baselineskip}
\begin{practiceans}
  \correct{D}
  Rearrange: $T^2 = \dfrac{4\pi^2}{GM}\, r^3$. This has the form $y = kx$ with $y = T^2$ and
  $x = r^3$. So $T^2$ against $r^3$ is a straight line through the origin. Its gradient is
  $\dfrac{4\pi^2}{GM}$.
  \begin{itemize}
    \item \xmark\ A, B and C: each of these pairs gives a curve. For example $T \propto r^{3/2}$, so $T$
          against $r$ curves upwards.
  \end{itemize}
\end{practiceans}

% ---- Q13
\Needspace{9\baselineskip}
\begin{practiceans}
  \correct{A}
  Kepler's second law: the line from the star to the planet sweeps out equal areas in equal times.
  The areas are equal, so the times are equal: 40 days. The arc from $P_1$ to $P_2$ is much longer
  than the arc from $P_3$ to $P_4$, but the planet covers it in the same time. So the planet moves
  faster near the star.
  \begin{itemize}
    \item \xmark\ B: the planet is fastest near the star, not furthest from it.
    \item \xmark\ C and D: equal areas \emph{always} mean equal times.
  \end{itemize}
\end{practiceans}

% ---- Q14
\Needspace{9\baselineskip}
\begin{practiceans}
  \correct{C}
  For everything that orbits the Sun, $\dfrac{r^3}{T^2}$ is the same. For Earth, $r = \qty{1}{AU}$ and
  $T = \qty{1}{\yr}$, so $\dfrac{r^3}{T^2} = 1$. For the dwarf planet:
  \[ r^3 = T^2 = 125^2 = 15\,625 \quad\Rightarrow\quad r = \sqrt[3]{15\,625} = \qty{25}{AU} \]
  \begin{itemize}
    \item \xmark\ A (\qty{5.0}{AU}): uses $r^3 = T$, the cube root of 125.
    \item \xmark\ B (\qty{11}{AU}): uses $r^2 = T$, the square root of 125.
    \item \xmark\ D (\qty{1400}{AU}): swaps $r$ and $T$: $r^2 = T^3$.
  \end{itemize}
\end{practiceans}

% ---- Q15
\Needspace{9\baselineskip}
\begin{practiceans}
  \correct{A}
  The radius is measured from the centre of Earth:
  $r = \num{6.371e6} + \num{4.20e5} = \qty{6.791e6}{m}$.
  \[ v = \sqrt{\frac{GM}{r}} = \sqrt{\frac{(\num{6.67e-11})(\num{6.0e24})}{\num{6.791e6}}}
       = \qty{7.68e3}{m.s^{-1}} \]
  \begin{itemize}
    \item \xmark\ B: uses $r = r_\mathrm{E}$ and forgets the altitude.
    \item \xmark\ C: uses the escape velocity formula $\sqrt{\dfrac{2GM}{r}}$.
    \item \xmark\ D: uses only the altitude (\qty{4.20e5}{m}) as $r$.
  \end{itemize}
\end{practiceans}

% ---- Q16
\Needspace{9\baselineskip}
\begin{practiceans}
  \correct{B}
  Change the period to seconds: $T = 24.0 \times 3600 = \qty{86400}{s}$.
  \[ r = \sqrt[3]{\frac{GMT^2}{4\pi^2}}
       = \sqrt[3]{\frac{(\num{6.67e-11})(\num{6.0e24})(86\,400)^2}{4\pi^2}}
       = \qty{4.23e7}{m} \]
  This is measured from the centre of Earth. The altitude is
  $\num{4.23e7} - \num{6.371e6} = \qty{3.59e7}{m}$ (about \qty{36000}{km}).
  \begin{itemize}
    \item \xmark\ A: uses $T = 24.0$ (hours) instead of seconds. This radius is smaller than
          Earth, so it is impossible.
    \item \xmark\ C: this is the radius of the orbit. It forgets to subtract $r_\mathrm{E}$.
    \item \xmark\ D: forgets the cube root. This number is $r^3$.
  \end{itemize}
\end{practiceans}

% ---- Q17
\Needspace{9\baselineskip}
\begin{practiceans}
  \correct{C}
  Rearrange $\dfrac{r^3}{T^2} = \dfrac{GM}{4\pi^2}$ for $M$:
  \[ M = \frac{4\pi^2 r^3}{GT^2}
       = \frac{4\pi^2 (\num{4.22e8})^3}{(\num{6.67e-11})(\num{1.53e5})^2}
       = \qty{1.90e27}{kg} \]
  \begin{itemize}
    \item \xmark\ A: uses $r^2$ instead of $r^3$.
    \item \xmark\ B: forgets the $4\pi^2$.
    \item \xmark\ D: forgets to square $T$.
  \end{itemize}
\end{practiceans}

% ---- Q18
\Needspace{9\baselineskip}
\begin{practiceans}
  \correct{D}
  $T^2 = \dfrac{4\pi^2}{GM}\, r^3$, so the gradient of the line is $\dfrac{4\pi^2}{GM}$. The
  planet's mass $M$ is on the bottom of the fraction. A bigger mass gives a \emph{smaller}
  gradient. Planet Y's line is less steep, so Y has the greater mass.
  \begin{itemize}
    \item \xmark\ A: a steeper line means a \emph{smaller} mass.
    \item \xmark\ B: both lines pass through the origin, but their gradients are different, so the
          masses are different.
    \item \xmark\ C: the moons' masses cancel out. $\dfrac{r^3}{T^2} = \dfrac{GM}{4\pi^2}$ contains
          only the planet's mass.
  \end{itemize}
\end{practiceans}

% ===================================================================== Section C
% ---- Q19
\Needspace{9\baselineskip}
\begin{practiceans}
  \correct{A}
  $U = -\dfrac{GMm}{r}$. At $Q$, $r$ is largest, so $U$ is the least negative. A less negative
  number is a greater number, so $U$ is greatest at $Q$. The total energy stays the same, so when
  $U$ goes up, the kinetic energy goes down.
  \begin{itemize}
    \item \xmark\ B and C: kinetic energy and speed are smallest at $Q$. The comet is slowest when it
          is furthest from the star.
    \item \xmark\ D: only gravity acts, so the total energy is the same at every point of the orbit.
  \end{itemize}
\end{practiceans}

% ---- Q20
\Needspace{9\baselineskip}
\begin{practiceans}
  \correct{B}
  $X$ is closer to Earth, so $U$ is more negative at $X$. The total energy stays the same, so the
  kinetic energy is bigger at $X$: the satellite is faster at $X$. Only gravity does work on the
  satellite (there is no air resistance), so the total energy $K + U$ is the same at $X$ and $Y$.
  \begin{itemize}
    \item \xmark\ A and C: the total energy does not change around the orbit.
    \item \xmark\ C and D: the satellite is fastest at the closest point, $X$, not at $Y$.
  \end{itemize}
\end{practiceans}

% ---- Q21
\Needspace{9\baselineskip}
\begin{practiceans}
  \correct{C}
  GEO is much further from Earth than LEO ($r \approx \qty{4.2e7}{m}$ compared with about
  \qty{7e6}{m}).
  \begin{itemize}
    \item Potential energy: $U = -\dfrac{GMm}{r}$. A bigger $r$ makes $U$ less negative, so $U$ is
          greater in \textbf{GEO}.
    \item Speed: $v = \sqrt{\dfrac{GM}{r}}$. A smaller $r$ gives a bigger speed, so the speed is
          greater in \textbf{LEO}.
  \end{itemize}
  \begin{itemize}
    \item \xmark\ A and B: “more negative” means \emph{smaller}. LEO has the more negative
          $U$, so it has the smaller $U$.
    \item \xmark\ B and D: a higher orbit is slower.
  \end{itemize}
\end{practiceans}

% ---- Q22
\Needspace{9\baselineskip}
\begin{practiceans}
  \correct{D}
  The total energy is $E = -\dfrac{GMm}{2r}$. A bigger $r$ makes $E$ less negative, so the total
  energy increases. This extra energy comes from the work done by the rocket engine.
  \begin{itemize}
    \item \xmark\ A: $K = \dfrac{GMm}{2r}$. A higher orbit is slower, so $K$ decreases.
    \item \xmark\ B: $U = -\dfrac{GMm}{r}$ becomes less negative, so $U$ increases.
    \item \xmark\ C: $T^2 \propto r^3$, so the period increases.
  \end{itemize}
\end{practiceans}

% ---- Q23
\Needspace{9\baselineskip}
\begin{practiceans}
  \correct{B}
  $r = r_\mathrm{E} + \text{altitude} = \num{6.371e6} + \num{3.50e6} = \qty{9.871e6}{m}$.
  \[ U = -\frac{GMm}{r} = -\frac{(\num{6.67e-11})(\num{6.0e24})(\num{1.20e3})}{\num{9.871e6}}
       = \qty{-4.87e10}{J} \]
  \begin{itemize}
    \item \xmark\ A: this is the total energy $E = \frac{1}{2}U$, not $U$.
    \item \xmark\ C: uses $r = r_\mathrm{E}$ and forgets the altitude.
    \item \xmark\ D: uses only the altitude as $r$.
  \end{itemize}
\end{practiceans}

% ---- Q24
\Needspace{9\baselineskip}
\begin{practiceans}
  \correct{C}
  Read the graph: at \qty{2000}{km}, $U \approx \qty{-9.6e9}{J}$. At \qty{12000}{km},
  $U \approx \qty{-4.4e9}{J}$.
  \[ \Delta U = U_\mathrm{final} - U_\mathrm{initial} = (\num{-4.4e9}) - (\num{-9.6e9})
     = \qty[retain-explicit-plus]{+5.2e9}{J} \]
  The change is positive: moving away from Earth increases $U$, because work must be done against
  gravity.
  \begin{itemize}
    \item \xmark\ A: adds the two values: $(-9.6) + (-4.4) = -14$, so \qty{-1.4e10}{J}.
    \item \xmark\ B: the right size, but the wrong sign. It does initial $-$ final.
    \item \xmark\ D: adds the two sizes: $9.6 + 4.4 = 14$.
  \end{itemize}
\end{practiceans}

% ---- Q25
\Needspace{9\baselineskip}
\begin{practiceans}
  \correct{A}
  $F = \dfrac{GMm}{r^2}$. The force becomes $\frac{1}{9}$, so $r^2$ is 9 times bigger and $r$ is
  3 times bigger. $U = -\dfrac{GMm}{r}$, so $U \propto \dfrac{1}{r}$: $U$ becomes $\frac{1}{3}$ of
  its first value. For example, \qty{-9e9}{J} becomes \qty{-3e9}{J}. (This is an
  \emph{increase}, because the value is less negative.)
  \begin{itemize}
    \item \xmark\ B: uses $U \propto \frac{1}{r^2}$, like the force. $U$ has only $r$ on the bottom.
    \item \xmark\ C: multiplies $U$ by the change in $r$ instead of dividing.
    \item \xmark\ D: multiplies by the change in $r^2$ instead of dividing by the change in $r$.
  \end{itemize}
\end{practiceans}

% ---- Q26
\Needspace{9\baselineskip}
\begin{practiceans}
  \correct{C}
  $v_\mathrm{esc} = \sqrt{\dfrac{2GM}{r}}$. It depends on the planet's mass $M$ and the distance
  $r$. It does \emph{not} depend on the probe's mass. The distance goes from $8r$ to $2r$, so it
  is divided by 4. Then $v_\mathrm{esc}$ is multiplied by $\sqrt{4} = 2$. The answer is $2v$.
  \begin{itemize}
    \item \xmark\ A: the ratio is upside down. Closer to Earth, the escape velocity is bigger.
    \item \xmark\ B: uses the change in the probe's mass. The probe's mass does not matter.
    \item \xmark\ D: forgets the square root.
  \end{itemize}
\end{practiceans}

% ---- Q27
\Needspace{9\baselineskip}
\begin{practiceans}
  \correct{A}
  $v_\mathrm{esc} = \sqrt{\dfrac{2GM}{r}}$, where $r$ is the radius (half the diameter). Compare
  $\dfrac{M}{r}$ for each planet:
  \begin{itemize}
    \item X: $\dfrac{\num{2.0e24}}{\num{4.0e6}} = \num{5.0e17}$
          \quad(about \qty{8.2}{km.s^{-1}})
    \item Y: $\dfrac{\num{6.0e24}}{\num{1.0e7}} = \num{6.0e17}$
          \quad(about \qty{8.9}{km.s^{-1}})
    \item Z: $\dfrac{\num{1.0e24}}{\num{1.5e6}} = \num{6.7e17}$
          \quad(about \qty{9.4}{km.s^{-1}})
  \end{itemize}
  So Z is largest, then Y, then X. The distance from the star does not matter.
  \begin{itemize}
    \item \xmark\ B: ranks the planets by mass only.
    \item \xmark\ C: ranks the planets by size only (smallest first).
    \item \xmark\ D: ranks the planets by their distance from the star. This has nothing to do with
          escape velocity from the surface.
  \end{itemize}
\end{practiceans}

% ---- Q28
\Needspace{9\baselineskip}
\begin{practiceans}
  \correct{B}
  \[ v_\mathrm{esc} = \sqrt{\frac{2GM}{r}}
       = \sqrt{\frac{2(\num{6.67e-11})(\num{3.30e23})}{\num{2.44e6}}}
       = \qty{4.25e3}{m.s^{-1}} \]
  \begin{itemize}
    \item \xmark\ A: this is the orbital speed $\sqrt{\dfrac{GM}{r}}$. It forgets the 2.
    \item \xmark\ C: uses $r$ in kilometres instead of metres.
    \item \xmark\ D: forgets the square root.
  \end{itemize}
\end{practiceans}

% ---- Q29
\Needspace{9\baselineskip}
\begin{practiceans}
  \correct{D}
  Each planet adds its own (negative) potential energy:
  \[ U = -\frac{GMm}{r_\mathrm{P}} - \frac{GMm}{r_\mathrm{Q}} \]
  At $S$, the mass is very close to P, so $U$ is very negative (low). In the middle, the mass is
  far from both planets, so $U$ rises to its highest (least negative) value. Near $F$, it is close
  to Q, so $U$ falls again. $U$ rises, then falls, and stays below zero. The mass moves at a
  steady speed, so the time axis is like a position axis.
  \begin{itemize}
    \item \xmark\ A: the distances to the planets change, so $U$ changes.
    \item \xmark\ B: this shape has $U$ lowest in the middle. But the middle is where the mass is
          furthest from both planets.
    \item \xmark\ C: a steady rise ignores planet Q. Near $F$ the mass is close to Q, so $U$ must
          fall again.
  \end{itemize}
\end{practiceans}

% ---- Q30
\Needspace{9\baselineskip}
\begin{practiceans}
  \correct{B}
  The energy supplied is the change in total energy:
  \[ \Delta E = \left(-\frac{GMm}{2r_2}\right) - \left(-\frac{GMm}{2r_1}\right)
             = \frac{GMm}{2}\left(\frac{1}{r_1} - \frac{1}{r_2}\right) \]
  \[ \Delta E = \frac{(\num{6.67e-11})(\num{6.0e24})(800)}{2}
       \left(\frac{1}{\num{7.0e6}} - \frac{1}{\num{2.1e7}}\right) = \qty{1.5e10}{J} \]
  \begin{itemize}
    \item \xmark\ A: this is only the size of the total energy in the new orbit.
    \item \xmark\ C: this is the change in $U$ only. In the higher orbit the satellite is slower,
          so it loses kinetic energy. This gives back half of the increase in $U$.
    \item \xmark\ D: this is the size of $U$ in the first orbit.
  \end{itemize}
\end{practiceans}

\end{practiceanswers}

\end{document}
